\section{Taxonomy of LLM-System Quality Assessment}
\label{sec:taxonomy}

The survey adopts a two-layer taxonomy. The first layer represents \emph{what is evaluated}: the lifecycle dimensions of an LLM-based system. The second represents \emph{how quality is evaluated}: the methods and protocols used to produce evidence about those dimensions. Evaluator reliability is positioned outside the normal execution lifecycle as a meta-evaluation layer.

\subsection{Lifecycle Dimensions}

The six lifecycle dimensions are summarized below.

\begin{description}

\item[Input and Specification Quality:]
Evaluates whether the task is adequately specified and whether system behavior is robust to acceptable variations in prompts, instructions, formatting, constraints, or other input-level choices. Typical concerns include ambiguity, prompt sensitivity, instruction following, and specification robustness.

\item[Context and Retrieval Quality:]
Evaluates the external information supplied to the model, including retrieval, ranking, grounding, context selection, memory retrieval, and evidence coverage. Typical concerns include relevance, recall, precision, context sufficiency, grounding, and retrieval efficiency.

\item[Output Quality:]
Evaluates the generated response or artifact. Typical concerns include correctness, factuality, faithfulness, relevance, completeness, coherence, helpfulness, safety, code correctness, and task-specific performance.

\item[Conversation and Interaction Quality:]
Evaluates behavior across multiple turns or interactive sessions. Typical concerns include consistency, state maintenance, adaptation to user feedback, dialogue-policy adherence, long-horizon behavior, and user-goal completion.

\item[Agent, Process, and Trajectory Quality:]
Evaluates intermediate decisions and execution behavior in agentic systems. Typical concerns include planning, decomposition, routing, tool selection, tool correctness, intermediate reasoning, trajectory quality, coordination, and progress toward the task goal.

\item[Operational and System Quality:]
Evaluates whether the system can operate acceptably under deployment constraints. Typical concerns include latency, throughput, monetary cost, token consumption, memory and compute usage, energy consumption, availability, privacy, safety, and observability.

\end{description}

\subsection{Cross-Cutting Properties}

Robustness, safety and security, fairness and bias, multimodality, domain specificity, and cost/efficiency are treated as cross-cutting properties rather than lifecycle stages. A robustness study, for example, may concern prompts, retrieval, outputs, agents, or evaluators. Recording the lifecycle target and the cross-cutting property separately preserves this distinction.

\subsection{Evaluation-Method Families}

Evaluation methods are classified independently of the dimension they assess:

\begin{enumerate}
    \item deterministic and rule-based evaluation;
    \item lexical or reference-based metrics;
    \item semantic or reference-aware metrics;
    \item learned evaluators or classifiers;
    \item human evaluation;
    \item LLM-as-a-Judge;
    \item multi-judge or ensemble evaluation;
    \item task- or environment-based verification;
    \item adversarial or red-team evaluation;
    \item perturbation or mutation testing;
    \item trace-aware evaluation;
    \item statistical reliability or meta-evaluation;
    \item hybrid approaches.
\end{enumerate}

The same method can be used at different lifecycle stages. For example, an LLM judge may score output factuality, assess a plan, evaluate a tool-use trajectory, or compare alternative responses. Accordingly, method frequency is not interpreted as evidence that a particular quality dimension is well covered.

\subsection{Evaluator Reliability as a Meta-Evaluation Layer}

Evaluator reliability is treated separately because it describes the quality of the measurement instrument. Relevant properties include inter- and intra-evaluator consistency, agreement with humans, calibration, bias, prompt sensitivity, presentation effects, robustness to adversarial or persuasive inputs, and dependence among evaluators.

This distinction is particularly important for LLM-as-a-Judge. The statement that a study \emph{uses} an LLM judge does not imply that it evaluates the reliability of that judge. The latter requires explicit evidence about the judge itself.

\subsection{Conceptual Overview}

Figure~\ref{fig:taxonomy} will present the taxonomy as a circular lifecycle. The inner structure represents the six system-quality dimensions. Evaluation methods form a second layer that can act on any dimension. The outer layer represents evaluator reliability and other meta-evaluation concerns.

\begin{figure*}[t]
    \centering
    \fbox{
        \parbox{0.92\textwidth}{
        \centering
        \textbf{Placeholder for final circular taxonomy}\\[5pt]
        Lifecycle dimensions: Input/Specification $\rightarrow$ Context/Retrieval
        $\rightarrow$ Output $\rightarrow$ Conversation/Interaction
        $\rightarrow$ Agent/Process/Trajectory $\rightarrow$ Operational/System\\[5pt]
        Evaluation layer: rules, reference metrics, semantic metrics, learned evaluators,
        humans, LLM judges, multi-judge schemes, task/environment checks, adversarial
        evaluation, perturbation, trace-aware evaluation\\[5pt]
        Meta-evaluation layer: reliability, agreement, consistency, calibration,
        bias, robustness
        }
    }
    \caption{Conceptual organization of LLM-system quality assessment. The lifecycle defines what is evaluated; evaluation methods define how it is measured; evaluator reliability characterizes the measurement process itself.}
    \label{fig:taxonomy}
\end{figure*}

\begin{table*}[t]
\caption{Lifecycle dimensions and representative evaluation concerns.}
\label{tab:quality-dimensions}
\begin{tabular}{p{0.18\textwidth} p{0.29\textwidth} p{0.24\textwidth} p{0.21\textwidth}}
\toprule
\textbf{Dimension} & \textbf{Central question} & \textbf{Representative properties} & \textbf{Typical evidence} \\
\midrule
Input / Specification & Is the intended task adequately and robustly specified? & sensitivity, ambiguity, instruction quality, robustness & perturbation, task performance, human assessment \\
Context / Retrieval & Was useful and appropriate evidence supplied? & relevance, recall, grounding, coverage, context sufficiency & IR metrics, semantic evaluation, human/LLM assessment \\
Output / Response & Is the resulting artifact correct and useful? & correctness, factuality, faithfulness, relevance, safety & reference metrics, humans, judges, execution tests \\
Conversation / Interaction & Does quality persist across turns and feedback? & consistency, state, adaptation, goal completion & simulation, humans, task success, dialogue traces \\
Agent / Process / Trajectory & Was the task executed through an acceptable process? & planning, tools, intermediate decisions, coordination & traces, environment checks, judges, task outcomes \\
Operational / System & Can the system operate acceptably in practice? & latency, throughput, cost, availability, resources, privacy & telemetry, profiling, monitoring, load tests \\
Meta-evaluation & Can the quality judgment itself be trusted? & agreement, calibration, bias, consistency, robustness & human comparison, reliability statistics, perturbation \\
\bottomrule
\end{tabular}
\end{table*}
